% =========================== 1. Introduction ===========================
\section{Introduction}\label{sec:intro}

Attention has become a general computational primitive.  In machine
learning, content-based weighting over key--value pairs underlies
sequence transduction \citep{vaswani2017,bahdanau2015,luong2015},
pretrained language models \citep{devlin2019,brown2020} and is studied
mechanistically as a form of information routing
\citep{elhage2021,olsson2022}.  In neuroscience, attention-like gain and
selection have been modelled at the level of normalization
\citep{reynolds2009}, saliency \citep{koch1985,itti1998} and dynamic
routing of information \citep{olshausen1993}, while binding has classically
been associated with temporal correlation and synchrony
\citep{von_der_malsburg1995,von_der_malsburg1986,singer1995,treisman1980}.

Most of these mechanisms are studied at the population or network level.
This raises a more primitive question that is rarely asked: \emph{what
happens if attention-like temporal weighting is compressed into a single
dynamical unit?}  A single neuron is a scalar dynamical system with a
leaky membrane potential; classical integrate-and-fire dynamics
\citep{lapicque1907,abbott1999,hodgkin1952,burkitt2006,gerstner2014} have
no mechanism for query-dependent selection of stored content.  Whether a
minimal single-neuron architecture can nevertheless exhibit
history-dependent computation, and where its mechanistic boundary lies,
is an open question about the nature of attention as a computational
primitive rather than about a specific benchmark.

The Temporal Attention Neuron (TAN) provides a minimal laboratory for
this question.  TAN is a discrete-time single neuron that combines (i) a
temporal receptive field over the recent input stream, (ii) a
surprise-gated, Boltzmann-style weighting of that window -- an
attention-like saliency kernel -- and (iii) leaky spiking membrane
dynamics.  The presence of attention-like weighting inside a dynamical
unit does not, however, guarantee that the unit can \emph{route}
information: weighting strength is not content selection.

The purpose of this paper is therefore not to propose a better neuron,
but to perform a mechanistic decomposition of one minimal construction.
We ask, in sequence, four nested questions.  (i) \emph{State}: does
temporal history become part of the effective state, so that the scalar
membrane potential alone is insufficient?  (ii) \emph{Geometry}: does
attention-like weighting reorganize the local, event-locked response
geometry relative to uniform temporal integration?  (iii)
\emph{Computation}: does this history contribute to target-relevant
retrieval in a controlled temporal task?  (iv) \emph{Identifiability}:
does the scalar attention kernel actually implement query-conditioned
routing, i.e.\ a query determining which stored event is read out?

These capabilities form a hierarchy that must be verified layer by
layer:
\begin{equation*}
  \text{memory} \;\rightarrow\; \text{weighting} \;\rightarrow\;
  \text{routing} \;\rightarrow\; \text{binding}.
\end{equation*}
Each arrow is a separate computational claim.  The experiments reported
here establish the first two layers for TAN, provide limited evidence on
the third, and show that the fourth is not implemented by the current
scalar kernel.  The central thesis is:

\begin{quote}
\emph{A minimal temporal-attention neuron can acquire history-dependent,
non-Markovian dynamics and can reorganize event-locked response geometry
through surprise-gated temporal weighting; however, its scalar attention
kernel does not provide genuine query-conditioned content-addressable
routing.  Temporal memory, attention-like saliency weighting, routing
and semantic binding are therefore distinct computational capabilities,
not interchangeable consequences of increasing representational
complexity.}
\end{quote}

The remainder of the paper is organized as follows.
Section~\ref{sec:background} positions TAN with respect to state
sufficiency, input-dependent weighting, routing and mechanistic
identifiability.  Section~\ref{sec:model} gives the frozen model
definition.  Section~\ref{sec:methods} describes the experimental
program (state-sufficiency collisions, event-locked geometry, a temporal
distractor task, and an identifiability audit).  Section~\ref{sec:results}
reports results; Section~\ref{sec:discussion} interprets them;
Section~\ref{sec:limitations} lists limitations;
Section~\ref{sec:conclusion} concludes with the conceptual separation
that this paper establishes.
